\documentclass[a4paper,12pt]{article}

\usepackage[utf8]{inputenc}
\usepackage[T1]{fontenc}
\usepackage[slovak]{babel}
\usepackage{graphicx}
\usepackage{amsmath}
\usepackage{hyperref}
\usepackage{geometry}
\usepackage{tikz}
\usetikzlibrary{arrows.meta, positioning, babel}

\geometry{margin=2.5cm}

\title{Supervízne riadenie autonómnej skupiny dronov vo virtuálnej realite}
\author{Projektová dokumentácia}
\date{\today}

\begin{document}

%------------------------------------------------
% TITULNA STRANA
%------------------------------------------------
\begin{titlepage}
\centering

\vspace*{3cm}

{\LARGE \textbf{Supervízne riadenie autonómnych dronov pomocou virtuálnej reality}}

\vspace{1.5cm}

{\Large Projektová dokumentácia}

\vspace{2cm}

\textbf{Autori:} Bc. Marek Malina, Bc. Tadeáš Kapina, Bc. Šimon Maco, Bc. Miroslav Beňušík, Bc. Ľuboš Valkovič, Bc. Boris Šupák, Bc. Martin Brandobur, Bc. Ivan Vanca

\vfill

\today

\end{titlepage}

%------------------------------------------------
% OBSAH
%------------------------------------------------
\tableofcontents
\newpage

%------------------------------------------------
\section{Analýza zadania a požiadaviek}

Cieľom projektu je návrh a implementácia systému umožňujúceho supervízne riadenie skupiny autonómnych dronov pomocou virtuálnej reality (VR).
Drony vykonávajú autonómne inšpekčné misie vo vnútornom priestore, pričom operátor má možnosť sledovať ich stav a v prípade potreby
prevziať riadenie niektorého z dronov pomocou VR ovládačov.

Na rozdiel od jednoduchých systémov určených pre jeden dron musí navrhované riešenie podporovať viacero dronov súčasne.
Každý dron vykonáva vlastnú autonómnu misiu, pričom operátor má možnosť sledovať ich činnosť a vybrať konkrétny dron,
do ktorého riadenia môže zasiahnuť. Tento prístup umožňuje efektívne monitorovanie a koordináciu viacerých autonómnych
robotických systémov v rovnakom prostredí.

Hlavnou požiadavkou systému je bezpečný prechod medzi autonómnym a manuálnym režimom riadenia.
Operátor musí byť schopný zasiahnuť do riadenia bez destabilizácie letu a následne umožniť dronu
vrátiť sa späť do autonómneho režimu.
\\

\textbf{Hlavné požiadavky systému:}

\begin{itemize}
    \item autonómne vykonávanie misií viacerými dronmi
    \item možnosť výberu konkrétneho drona zo zoznamu dostupných zariadení
    \item zobrazenie telemetrických dát a video streamu z jednotlivých dronov
    \item manuálne riadenie vybraného drona pomocou VR ovládačov
    \item bezpečný návrat do autonómneho režimu po manuálnom zásahu operátora
\end{itemize}

Systém musí okrem zabezpečenia požadovanej funkcionality spĺňať aj viaceré prevádzkové a technické požiadavky.
Dôležitá je najmä spoľahlivá komunikácia medzi jednotlivými komponentmi systému,
nízka latencia pri prenose riadiacich príkazov a implementácia bezpečnostných mechanizmov,
ktoré minimalizujú riziko straty kontroly nad dronmi počas prevádzky.

%------------------------------------------------
\section{Návrh základnej architektúry systému}

Celý systém je navrhnutý ako distribuovaná architektúra pozostávajúca z troch hlavných častí:

\begin{itemize}
\item autonómny systém dronov
\item supervízna riadiaca vrstva
\item používateľské rozhranie vo virtuálnej realite
\end{itemize}

Autonómny systém každého drona zabezpečuje stabilizáciu letu, lokalizáciu a vykonávanie plánovanej misie.
Supervízna vrstva koordinuje jednotlivé drony a umožňuje operátorovi sledovať ich stav.
Používateľské rozhranie vo VR poskytuje operátorovi prehľad o všetkých aktívnych dronoch a
umožňuje zasiahnuť do riadenia vybraného zariadenia.

Komunikácia medzi jednotlivými komponentmi systému je realizovaná pomocou middleware ROS~2,
ktorý umožňuje modulárnu architektúru a jednoduchú integráciu jednotlivých komponentov.

\begin{figure}[h!]
\centering

\begin{tikzpicture}[
block/.style={draw, rectangle, rounded corners, align=center, minimum width=6cm, minimum height=1.4cm},
line/.style={-Latex, thick},
node distance=2cm
]

\node[block] (vr) {Používateľské rozhranie\\vo virtuálnej realite};

\node[block, below=of vr] (supervisor) {Supervízna riadiaca vrstva};

\node[block, below=of supervisor] (drone) {Autonómny systém dronov};

% komunikácia VR <-> supervisor
\draw[line] (vr) -- node[right]{riadiace príkazy} (supervisor);
\draw[line] (supervisor) -- node[left]{telemetria, stav systému} (vr);

% komunikácia supervisor <-> drony
\draw[line] (supervisor) -- node[right]{príkazy misie} (drone);
\draw[line] (drone) -- node[left]{telemetria, video stream} (supervisor);

\end{tikzpicture}

\caption{Bloková schéma architektúry navrhovaného systému}
\label{fig:system_architecture}

\end{figure}

Na obrázku \ref{fig:system_architecture} je znázornená základná architektúra navrhovaného systému.
Používateľ komunikuje so systémom prostredníctvom rozhrania vo virtuálnej realite, ktoré umožňuje
vizualizáciu telemetrických údajov a manuálne riadenie vybraného drona. Supervízna riadiaca vrstva
zabezpečuje dohľad nad jednotlivými dronmi, výber aktívneho drona a prepínanie medzi autonómnym
a manuálnym režimom riadenia. Autonómny systém dronov vykonáva samotnú misiu a zabezpečuje
stabilizáciu letu, lokalizáciu a navigáciu jednotlivých dronov. Komunikácia medzi jednotlivými
vrstvami systému zahŕňa prenos riadiacich príkazov smerom k dronom a prenos telemetrie a video
streamu smerom k operátorovi.

%------------------------------------------------
\section{Zvolené hybridné riešenie}

Pre realizáciu projektu bolo zvolené hybridné riešenie, ktoré kombinuje simuláciu viacerých dronov
a experimentálne overenie na reálnom zariadení. Tento prístup umožňuje splniť požiadavku na
supervízne riadenie skupiny dronov a zároveň znižuje riziko spojené s testovaním viacerých
reálnych lietajúcich platforiem v jednom priestore.

Navrhované riešenie pozostáva z dvoch prevádzkových režimov:

\begin{itemize}
\item \textbf{simulačný režim} -- overenie správania systému na skupine troch dronov v prostredí Gazebo a ArduPilot SITL,
\item \textbf{reálny režim} -- overenie rovnakého princípu riadenia na jednom fyzickom drone.
\end{itemize}

V simulačnom režime sú spustené tri samostatné inštancie dronov, pričom každý dron má vlastný
MAVROS namespace, vlastnú autonómnu misiu a samostatný stavový automat. Operátor má možnosť
sledovať všetky drony, vybrať aktívny dron a dočasne prevziať jeho riadenie. Ostatné drony počas
tohto zásahu pokračujú vo svojich autonómnych misiách.

V reálnom režime sa používa jeden fyzický dron s rovnakým komunikačným rozhraním ako v simulácii.
Vďaka tomu je možné väčšinu softvérových komponentov použiť bez zmeny. Rozdiel je najmä v tom,
že namiesto simulovaného autopilota sa systém pripája na reálnu riadiacu jednotku drona cez
MAVLink a MAVROS.

\subsection{Hardvérové prepojenie systému}

Hardvérová architektúra systému je navrhnutá tak, aby bolo možné použiť rovnakú supervíznu vrstvu
pre simuláciu aj pre reálne zariadenie. V simulačnom režime bežia všetky komponenty na jednom
výkonnom počítači. V reálnom režime je operátorský počítač prepojený s dronom cez WiFi alebo lokálnu
sieť.

\begin{figure}[h!]
\centering

\begin{tikzpicture}[
block/.style={draw, rectangle, rounded corners, align=center, minimum width=4.8cm, minimum height=1.1cm},
smallblock/.style={draw, rectangle, rounded corners, align=center, minimum width=4.2cm, minimum height=1.0cm},
line/.style={-Latex, thick},
node distance=1.35cm
]

\node[block] (vr) {VR headset a ovládače};
\node[block, below=of vr] (pc) {Operátorský počítač\\ROS 2, supervízor, UI/VR};
\node[smallblock, below left=1.5cm and 0.7cm of pc] (sim) {Simulácia\\Gazebo + ArduPilot SITL};
\node[smallblock, below right=1.5cm and 0.7cm of pc] (real) {Reálny dron\\companion computer};
\node[smallblock, below=of real] (fc) {Riadiaca jednotka\\autopilot};

\draw[line] (vr) -- node[right]{vstupy operátora} (pc);
\draw[line] (pc) -- node[left]{ROS 2 / MAVROS} (sim);
\draw[line] (pc) -- node[right]{WiFi / LAN} (real);
\draw[line] (real) -- node[right]{MAVLink} (fc);
\draw[line] (real.west) to[out=180,in=0] node[left, align=center]{video stream\\telemetria} (pc.east);

\end{tikzpicture}

\caption{Hardvérové prepojenie hybridného riešenia}
\label{fig:hybrid_hardware}

\end{figure}

Operátorský počítač zabezpečuje beh supervíznej logiky, používateľského rozhrania, príjem
telemetrie a príjem video streamu. VR headset je pripojený k operátorskému počítaču a slúži ako
rozhranie pre výber drona, sledovanie stavu a manuálne riadenie. Reálny dron obsahuje companion
computer, ktorý sprostredkuje komunikáciu medzi sieťou, kamerou a riadiacou jednotkou drona.

Takto navrhnuté prepojenie umožňuje najprv bezpečne otestovať správanie systému v simulácii a až
následne preniesť overené časti riešenia na reálnu platformu.

%------------------------------------------------
\section{Návrh komunikačných rozhraní a ROS 2 uzlov}

Systém je implementovaný pomocou architektúry ROS~2 založenej na uzloch (nodes), ktoré komunikujú pomocou tém (topics), služieb (services) a akcií (actions).

Základné uzly systému sú:

\begin{itemize}
\item \textbf{Flight Interface Node} – komunikácia medzi ROS~2 a riadiacou jednotkou drona,
\item \textbf{State Estimation Node} – výpočet polohy a orientácie drona,
\item \textbf{Mission Executor Node} – vykonávanie plánovanej trajektórie,
\item \textbf{Supervisor Node} – riadenie režimov systému,
\item \textbf{VR Interface Node} – komunikácia medzi VR aplikáciou a ROS~2.
\end{itemize}

Medzi hlavné komunikačné témy patria napríklad:

\begin{itemize}
\item \texttt{/droneX/state} -- stav autopilota, režim letu a informácia o pripojení,
\item \texttt{/droneX/local\_position/pose} -- aktuálna poloha a orientácia drona,
\item \texttt{/droneX/battery} -- stav batérie,
\item \texttt{/droneX/setpoint\_position/local} -- pozičné setpointy autonómnej misie,
\item \texttt{/droneX/setpoint\_velocity/cmd\_vel\_unstamped} -- rýchlostné príkazy pri manuálnom riadení,
\item \texttt{/supervisor/active\_drone} -- informácia o aktuálne vybranom drone,
\item \texttt{/supervisor/mode} -- stav supervízneho riadenia.
\end{itemize}

Symbol \texttt{droneX} predstavuje konkrétny namespace drona, napríklad \texttt{drone1},
\texttt{drone2} alebo \texttt{drone3}. Takáto architektúra umožňuje jednoduché rozšírenie systému
o ďalšie drony alebo nové funkcionality. V reálnom režime môže byť aktívny iba jeden namespace,
pričom zvyšok softvérovej architektúry zostáva zachovaný.

Supervízny uzol prijíma požiadavky z používateľského rozhrania, vyhodnocuje ich a následne posiela
riadiace príkazy iba vybranému dronu. Používateľské rozhranie teda nekomunikuje priamo s autopilotom,
ale cez supervíznu vrstvu, ktorá zabezpečuje kontrolu režimov a bezpečnostné obmedzenia.

%------------------------------------------------
\section{Návrh stavového automatu a handover mechanizmu}

Riadenie režimov systému je realizované pomocou stavového automatu (Finite State Machine – FSM).
Tento automat zabezpečuje bezpečný prechod medzi jednotlivými režimami systému.

Základné stavy systému sú:

\begin{itemize}
\item \textbf{Idle} – systém čaká na spustenie misie,
\item \textbf{Autonomous Mission} – dron vykonáva autonómnu trajektóriu,
\item \textbf{Manual Control} – operátor riadi dron pomocou VR ovládačov,
\item \textbf{Hold} – dron udržiava stabilnú polohu,
\item \textbf{Landing} – bezpečné pristátie.
\end{itemize}

Pri prechode z autonómneho do manuálneho režimu systém najprv aktivuje stabilizačný režim
(\textit{Hold}), čím sa zabezpečí plynulý prechod bez náhlych zmien trajektórie.

Podobne pri návrate z manuálneho riadenia do autonómneho režimu systém vypočíta najbližší bod
trajektórie, od ktorého je možné bezpečne pokračovať v misii.

V prípade skupiny dronov sa handover vykonáva iba nad aktuálne vybraným dronom. Supervízor preto
udržiava premennú aktívneho drona a všetky manuálne príkazy smeruje výhradne na tento namespace.
Ostatné drony zostávajú v stave \textit{Autonomous Mission} a pokračujú vo vykonávaní vlastných
misií.

Typický priebeh prevzatia riadenia je nasledovný:

\begin{enumerate}
\item operátor vyberie konkrétny dron zo zoznamu dostupných zariadení,
\item používateľské rozhranie odošle požiadavku na prevzatie riadenia,
\item supervízor zastaví posielanie autonómnych setpointov vybranému dronu,
\item dron prejde do režimu \textit{Hold} a stabilizuje svoju polohu,
\item po stabilizácii supervízor povolí manuálne rýchlostné príkazy,
\item pri ukončení manuálneho režimu sa rýchlostné príkazy vynulujú a dron sa vráti do autonómnej misie.
\end{enumerate}

Pre prvú verziu riešenia je vhodné použiť jednoduchý mechanizmus návratu do misie. Dron si počas
autonómneho letu uchováva index aktuálneho waypointu. Po ukončení manuálneho zásahu pokračuje buď
od tohto waypointu, alebo od najbližšieho ešte nenavštíveného waypointu. Tento prístup je jednoduchý
na implementáciu a postačuje pre demonštračný scenár.

%------------------------------------------------
\section{Autonómna navigácia a vykonávanie misie}

Autonómna navigácia drona je založená na vykonávaní preddefinovanej trajektórie pozostávajúcej z
niekoľkých waypointov. Každý waypoint predstavuje cieľovú polohu, ktorú má dron dosiahnuť.

Počas letu systém priebežne vyhodnocuje aktuálnu polohu drona a generuje riadiace príkazy,
ktoré umožňujú sledovanie trajektórie.

Pre lokalizáciu vo vnútornom priestore je možné použiť rôzne metódy, napríklad:

\begin{itemize}
\item vizuálnu odometriu,
\item SLAM algoritmy,
\item externé lokalizačné systémy.
\end{itemize}

%------------------------------------------------
\section{VR používateľské rozhranie}

Používateľské rozhranie vo virtuálnej realite umožňuje operátorovi interagovať so systémom
intuitívnym spôsobom. Operátor má k dispozícii vizuálne zobrazenie prostredia, telemetrické údaje
a video stream z kamery drona.

Medzi hlavné funkcie VR rozhrania patria:

\begin{itemize}
\item výber dostupného drona zo zoznamu zariadení,
\item zobrazenie aktuálneho stavu drona,
\item vizualizácia video streamu,
\item manuálne riadenie pomocou VR ovládačov,
\item prepínanie medzi autonómnym a manuálnym režimom.
\end{itemize}

%------------------------------------------------
\section{Bezpečnostné mechanizmy a failsafe}

Bezpečnosť systému je kľúčovým aspektom návrhu. Systém musí byť schopný reagovať na rôzne
poruchové stavy, ako napríklad strata komunikácie alebo neplatné riadiace príkazy.

Medzi hlavné bezpečnostné mechanizmy patria:

\begin{itemize}
\item watchdog kontrolujúci dostupnosť komunikácie,
\item automatický prechod do stabilizačného režimu pri strate spojenia,
\item obmedzenie maximálnej rýchlosti a zrýchlenia,
\item núdzové pristátie pri kritickom stave batérie.
\end{itemize}



%------------------------------------------------
\section{Implementácia na reálnom zariadení}

Implementácia systému bude prebiehať postupne. Najskôr sa overí kompletný tok dát a riadiaca logika
v simulácii, následne sa rovnaká architektúra použije pri teste s reálnym dronom. Takýto postup
znižuje riziko poškodenia hardvéru a umožňuje jednoduchšie ladenie jednotlivých komponentov.

\subsection{Simulačný režim}

V simulačnom režime je použitý jeden operátorský počítač, na ktorom sú spustené nasledujúce
komponenty:

\begin{itemize}
\item Gazebo svet reprezentujúci vnútorné inšpekčné prostredie,
\item tri inštancie ArduPilot SITL,
\item tri inštancie MAVROS s namespacami \texttt{/drone1}, \texttt{/drone2} a \texttt{/drone3},
\item ROS~2 uzol pre vykonávanie autonómnych misií,
\item supervízny uzol pre výber aktívneho drona a handover mechanizmus,
\item používateľské rozhranie pre sledovanie telemetrie, video streamu a manuálne ovládanie.
\end{itemize}

Každý simulovaný dron má vlastnú misiu uloženú v samostatnom súbore s waypointmi. Misie sú navrhnuté
tak, aby drony vykonávali inšpekciu rôznych častí priestoru. Supervízor priebežne sleduje ich stav
a umožňuje operátorovi vybrať konkrétny dron na manuálny zásah.

\subsection{Reálny režim}

V reálnom režime sa použije jeden fyzický dron. Dron je vybavený riadiacou jednotkou, companion
computerom a kamerou. Companion computer zabezpečuje komunikáciu s operátorským počítačom, spustenie
MAVLink routra a prenos video streamu.

Hardvérové prepojenie reálneho režimu pozostáva z nasledujúcich častí:

\begin{itemize}
\item operátorský počítač so systémom ROS~2, MAVROS a používateľským rozhraním,
\item WiFi alebo lokálna sieť medzi počítačom a dronom,
\item companion computer na drone,
\item riadiaca jednotka drona komunikujúca cez MAVLink,
\item kamera posielajúca video stream na operátorský počítač.
\end{itemize}

V tomto režime sa overuje rovnaký princíp ako v simulácii, ale iba na jednom drone. Operátor spustí
autonómnu misiu, sleduje stav drona, následne prevezme manuálne riadenie a po krátkom zásahu vráti
dron späť do autonómneho režimu.

\subsection{Postup integrácie}

Implementácia bude rozdelená do niekoľkých krokov:

\begin{enumerate}
\item spustenie a overenie autonómnych misií troch dronov v simulácii,
\item doplnenie supervízneho uzla pre výber aktívneho drona,
\item implementácia handover mechanizmu medzi autonómnym a manuálnym riadením,
\item napojenie používateľského rozhrania a VR vstupov,
\item prenos overenej logiky na jeden reálny dron,
\item experimentálne meranie odozvy, stability a spoľahlivosti systému.
\end{enumerate}

%------------------------------------------------
\section{Experimentálne overenie}

Experimentálne overenie bude realizované v dvoch fázach. Prvá fáza prebehne v simulácii, kde je
možné bezpečne testovať aj hraničné situácie. Druhá fáza prebehne na reálnom drone v kontrolovaných
podmienkach.

\subsection{Overenie v simulácii}

V simulácii budú testované najmä nasledujúce scenáre:

\begin{itemize}
\item autonómny let troch dronov po samostatných trajektóriách,
\item zobrazenie telemetrie všetkých dronov v používateľskom rozhraní,
\item výber aktívneho drona operátorom,
\item prevzatie manuálneho riadenia vybraného drona počas autonómnej misie,
\item návrat vybraného drona do autonómneho režimu,
\item pokračovanie ostatných dronov v autonómnej misii počas manuálneho zásahu,
\item reakcia systému na výpadok manuálnych riadiacich príkazov.
\end{itemize}

\subsection{Overenie na reálnom drone}

Na reálnom zariadení bude cieľom overiť funkčnosť základného toku riadenia:

\begin{enumerate}
\item pripojenie operátorského počítača k dronu,
\item príjem telemetrie a video streamu,
\item spustenie autonómnej misie,
\item aktivácia režimu \textit{Hold},
\item manuálne riadenie pomocou ovládača alebo VR vstupu,
\item bezpečný návrat do autonómneho režimu,
\item bezpečné ukončenie letu.
\end{enumerate}

Počas experimentov budú sledované najmä tieto veličiny:

\begin{itemize}
\item latencia medzi vstupom operátora a reakciou drona,
\item čas potrebný na prechod z autonómneho do manuálneho režimu,
\item čas potrebný na návrat do autonómnej misie,
\item stabilita polohy počas režimu \textit{Hold},
\item spoľahlivosť prenosu telemetrie a video streamu.
\end{itemize}

%------------------------------------------------
\section{Diskusia obmedzení a ďalší rozvoj}

Navrhované hybridné riešenie predstavuje praktický kompromis medzi požiadavkou na riadenie skupiny
dronov a možnosťami bezpečného experimentovania s reálnym hardvérom. Simulácia umožňuje overiť
správanie viacerých dronov naraz, zatiaľ čo reálny test potvrdzuje použiteľnosť architektúry na
fyzickej platforme.

Hlavným obmedzením riešenia je rozdiel medzi simulovaným a reálnym prostredím. V simulácii sú
podmienky lepšie kontrolovateľné, zatiaľ čo reálny dron je ovplyvnený oneskorením komunikácie,
šumom senzorov, kvalitou lokalizácie, stavom batérie a fyzikálnymi obmedzeniami platformy. Z toho
dôvodu je potrebné všetky bezpečnostné parametre pred reálnym letom nastaviť konzervatívne.

Ďalším obmedzením je integrácia video streamu do VR prostredia. Prenos obrazu môže mať vyššiu
latenciu ako prenos riadiacich príkazov, preto musí byť manuálne riadenie navrhnuté tak, aby nebolo
závislé výhradne od video obrazu. Operátor by mal mať k dispozícii aj základnú telemetriu, stav
režimu a informáciu o polohe drona.

V ďalšom rozvoji projektu je možné systém rozšíriť o:

\begin{itemize}
\item pokročilejšie plánovanie návratu do misie po manuálnom zásahu,
\item detekciu a vyhýbanie sa kolíziám medzi dronmi,
\item podporu viacerých reálnych dronov,
\item realistickejšie VR prostredie so zobrazením mapy priestoru,
\item záznam experimentov a následnú analýzu letových dát,
\item automatické vyhodnocovanie kvality inšpekčnej misie.
\end{itemize}

Pre účely tímového projektu je však prioritou vytvoriť stabilný demonštrátor, ktorý jasne ukáže
základný princíp zdieľaného supervízneho riadenia: autonómny let, výber drona, manuálny zásah
operátora a bezpečný návrat do autonómnej misie.

\end{document}
